\section{A common invariant limit under uniform pure deaths}
\label{sec:common-limit}

\begin{discussion}
  \item Motivate the uniform pure-death assumption as a source of
  damping along outgoing-terminal patch trails and relaxation of end
  factors.
  \item Explain the three-part proof decomposition: confinement up to
  time \(T\), the occurrence of a later successful interaction, and no
  later successful interaction.
  \item Explain why the theorem first proves uniform convergence on
  \Mstar\ and then extends linearly to \Mminus.
\end{discussion}

\begin{theorem}[Common invariant limit]
\label{thm:common-limit}
Let \(\Lambda\) have polynomial volume growth of exponent \(D\), and
let \(\cL\) satisfy \cref{ass:standing-locality} and have
\patchpositivity. Suppose that there is \(\varepsilon>0\) such that
\begin{equation}
  c_i^1(\eta)\ge\varepsilon
  \qquad
  (i\in\Lambda,\ \eta\in\{0,1\}^{\Lambda}).
  \label{eq:uniform-death}
\end{equation}
Then there is an \invariantmeasure\ \(\pi\) such that, for every local
function \(f\), some \(K_f<\infty\) satisfies
\begin{equation}
  \sup_{\nu\in\Mminus}
  \abs{\nu(P_tf)-\pi(f)}
  \le
  K_f(1+t)^D e^{-\varepsilon t/2}.
  \label{eq:common-limit-bound}
\end{equation}
For every \(A\Subset\Lambda\),
\begin{equation}
  \pi(\chi_A)
  =
  \E_A\left[
    \prod_{P\in\cP}C(P)\,
    \ind\bigl(\abs{\cP}<\infty\bigr)
  \right],
  \label{eq:limit-moment-formula}
\end{equation}
where \(\cP\) is the full patch family of the dual started from
\((A,+)\).
\end{theorem}

\begin{corollary}[\uniformexponentialergodicity]
\label{cor:uniform-ergodicity}
Under the hypotheses of \cref{thm:common-limit}, if
\(\pstar\le\frac12\mathbf1\), then \(\pi\) is the unique
\invariantmeasure\ and
\[
  \sup_{\xi\in\{0,1\}^{\Lambda}}
  \abs{P_tf(\xi)-\pi(f)}
  \le
  K_f(1+t)^D e^{-\varepsilon t/2}
\]
for every local \(f\). In particular, the system has
\uniformexponentialergodicity\ in the sense of
\cref{def:uniform-exponential-ergodicity}.
\end{corollary}

\begin{remark}[Finite centered perturbations]
\label{rem:finite-centered-perturbations}
The conclusion of \cref{thm:common-limit} holds whenever, for some
\(K<\infty\),
\begin{equation}
  \nu(\chi_A^\star)
  \ge
  -K\chi_A^\star(\mathbf1)
  \qquad(A\Subset\Lambda).
  \label{eq:finite-centered-perturbation}
\end{equation}
Indeed,
\[
  \frac{\nu+K\mu_{\mathbf1}}{1+K}\in\Mstar.
\]
The constant increases by at most \(1+2K\). Every deterministic
configuration with finitely many zeroes satisfies
\eqref{eq:finite-centered-perturbation}.
\end{remark}

\subsection{Less-noisy domination}

\begin{notation}[Less-noisy system and skeleton weights]
\label{not:less-noisy-weights}
Define \(\widehat{\cL}\) by
\begin{equation}
  \cL
  =
  \widehat{\cL}
  +
  \varepsilon\sum_{i\in\Lambda}
  \bigl(f(\eta^{i,0})-f(\eta)\bigr).
  \label{eq:less-noisy-generator}
\end{equation}
Let \(\widehat C\) be its patch factors. For
\(\nu\in\Mstar\), define
\begin{equation}
  W_t^\nu
  =
  \prod_{P\in\cB_t}C(P)\,
  \nu\left[
    \prod_{P\in\cE_t}
    C(\eta(i(P)),P)
  \right],
  \label{eq:finite-weight}
\end{equation}
and define \(\widehat W_t^\nu\) analogously. Finally, put
\begin{equation}
  W
  =
  \prod_{P\in\cP}C(P)\,
  \ind\bigl(\abs{\cP}<\infty\bigr).
  \label{eq:limiting-weight-W}
\end{equation}
\end{notation}

\begin{lemma}[Less-noisy domination]
\label{lem:less-noisy-domination}
The generator \(\widehat{\cL}\) has \patchpositivity\ and critical
profile \pstar. For every full patch,
\begin{equation}
  0\le C(P)\le\widehat C(P).
  \label{eq:full-factor-comparison}
\end{equation}
If \(P\) has outgoing terminal type, then
\begin{equation}
  C(P)
  =
  e^{-\varepsilon\Delta(P)}
  \widehat C(P).
  \label{eq:exact-XO-damping}
\end{equation}
Moreover,
\begin{equation}
  0\le W_t^\nu\le\widehat W_t^\nu,
  \qquad
  \E_A[W_t^\nu]=\nu(P_t\chi_A),
  \qquad
  \E_A[\widehat W_t^\nu]\le1.
  \label{eq:weight-domination}
\end{equation}
\end{lemma}

\begin{proof}
\Cref{thm:pure-death-comparison} and its coefficientwise proof give
\eqref{eq:full-factor-comparison} and the corresponding comparison of
centered end values and slopes. For an outgoing-terminal patch, the
source remains active throughout its lifetime, and the removed death
contributes exactly the factor in
\eqref{eq:exact-XO-damping}. Expanding the end product as in
\eqref{eq:end-factor-expansion} gives the first inequality in
\eqref{eq:weight-domination}. The equality follows from
\cref{thm:patch-representation}, while
\[
  \E_A[\widehat W_t^\nu]=\nu(\widehat P_t\chi_A)\le1.
\]
\end{proof}

\subsection{Confinement and late interactions}

\begin{notation}
\label{not:confinement-events}
Fix \(0\le T<t\) and a finite \(R\supseteq A\). Let
\[
  E_T^R=\{\Cone_T\subseteq R\},
\qquad
  \rho_A(T,R)
  =
  \norm{(P_T-P_T^{R,0})\chi_A}_\infty,
\]
where \(P_T^{R,0}\) is the zero-boundary semigroup on \(R\). Also set
\begin{align*}
  L_{T,t}
  &=
  \{\text{no successful interaction occurs in }(T,t]\},
  \\
  L_T
  &=
  \{\text{no successful interaction occurs after }T\}.
\end{align*}
\end{notation}

\begin{lemma}[Confinement bounds]
\label{lem:confinement-bounds}
For every \(\nu\in\Mstar\),
\begin{align}
  0
  &\le
  \E_A\left[
    W_t^\nu\ind((E_T^R)^c)
  \right]
  \le
  \rho_A(T,R),
  \label{eq:finite-confinement-error}
  \\
  0
  &\le
  \E_A\left[
    W\ind((E_T^R)^c)
  \right]
  \le
  \rho_A(T,R).
  \label{eq:limit-confinement-error}
\end{align}
\end{lemma}

\begin{proof}
\Cref{prop:confined-duality} gives
\[
  \E_A\left[W_t^\nu\ind(E_T^R)\right]
  =
  \nu(P_{t-T}P_T^{R,0}\chi_A).
\]
Subtract this identity from
\(\E_A[W_t^\nu]=\nu(P_t\chi_A)\), and use positivity and contraction
to obtain \eqref{eq:finite-confinement-error}. Finite-horizon all-one
weights converge to \(W\) on \(\{\abs{\cP}<\infty\}\). Positivity and
Fatou's lemma give \eqref{eq:limit-confinement-error}.
\end{proof}

\begin{lemma}[Late-interaction bounds]
\label{lem:late-interaction-bounds}
For every \(\nu\in\Mstar\),
\begin{align}
  0
  &\le
  \E_A\left[
    W_t^\nu\ind(E_T^R\cap L_{T,t}^c)
  \right]
  \le
  e^{-\varepsilon T},
  \label{eq:finite-late-bound}
  \\
  0
  &\le
  \E_A\left[
    W\ind(E_T^R\cap L_T^c)
  \right]
  \le
  e^{-\varepsilon T}.
  \label{eq:limit-late-bound}
\end{align}
\end{lemma}

\begin{proof}
On \(L_{T,t}^c\), follow the first successful interaction after \(T\)
backward through the outgoing-terminal patches leading to time zero.
This trail consists of bulk patches and has total lifetime at least
\(T\). Apply \eqref{eq:exact-XO-damping} on the trail,
\eqref{eq:full-factor-comparison} elsewhere, and then
\eqref{eq:weight-domination}. This proves
\eqref{eq:finite-late-bound}. The same argument for the full patch
product proves \eqref{eq:limit-late-bound}.
\end{proof}

\subsection{Relaxation without late interactions}

\begin{lemma}[No-late-interaction relaxation]
\label{lem:no-late-relaxation}
For \(Q\in\cE_T\), let \(Q^{\uparrow t}\) be its continuation through
time \(t\) without a successful interaction, including the conditional
continuation probability in its factor. Then
\begin{equation}
  \E_A\left[
    W_t^\nu\ind(L_{T,t})
    \,\middle|\,\cG_T
  \right]
  =
  \prod_{P\in\cB_T}C(P)\,
  \nu\left[
    \prod_{Q\in\cE_T}
    C(\eta(i(Q)),Q^{\uparrow t})
  \right],
  \label{eq:no-late-finite-identity}
\end{equation}
and
\begin{equation}
  \E_A\left[
    W\ind(L_T)
    \,\middle|\,\cG_T
  \right]
  =
  \prod_{P\in\cB_T}C(P)
  \prod_{Q\in\cE_T}C(Q^{\uparrow\infty}).
  \label{eq:no-late-limit-identity}
\end{equation}
Moreover,
\begin{equation}
  \abs{
    \E_A\left[
      W_t^\nu\ind(E_T^R\cap L_{T,t})
    \right]
    -
    \E_A\left[
      W\ind(E_T^R\cap L_T)
    \right]
  }
  \le
  (1+\abs R)e^{-\varepsilon(t-T)}.
  \label{eq:no-late-bound}
\end{equation}
\end{lemma}

\begin{proof}
On \(L_{T,t}\), each end patch at time \(T\) continues without a
successful interaction, which gives
\eqref{eq:no-late-finite-identity} by finite-horizon factorization.
\Cref{prop:continuation-relaxation} gives exponential relaxation with
rate
\[
  r_i=c_i^0(\vn)+c_i^1(\vn)\ge\varepsilon.
\]
Apply \eqref{eq:no-late-finite-identity} first at finite horizon \(u\)
and send \(u\to\infty\) to obtain
\eqref{eq:no-late-limit-identity}; no all-time factorization is used.

In the centered expansion of the end product, every nonconstant term
contains a slope and is bounded by
\(e^{-\varepsilon(t-T)}\) times the less-noisy weight. The constant
terms are bounded by telescoping the product, with one error per end
patch. On \(E_T^R\),
\(\abs{\cE_T}=\abs{\Cone_T}\le\abs R\). Integration and
\eqref{eq:weight-domination} yield \eqref{eq:no-late-bound}.
\end{proof}

\subsection{Proof of the common-limit theorem}

\begin{proof}[Proof of \cref{thm:common-limit}]
\Cref{lem:confinement-bounds,lem:late-interaction-bounds,lem:no-late-relaxation}
give, uniformly over \(\nu\in\Mstar\),
\begin{equation}
  \abs{\nu(P_t\chi_A)-\E_A[W]}
  \le
  2\rho_A(T,R)
  +
  2e^{-\varepsilon T}
  +
  (1+\abs R)e^{-\varepsilon(t-T)}.
  \label{eq:master-common-limit-bound}
\end{equation}
By \cref{lem:finite-propagation}, there is \(v<\infty\) such that, for
\(R_T=B(A,vT)\),
\[
  \rho_A(T,R_T)\le C_Ae^{-\varepsilon T},
  \qquad
  \abs{R_T}\le C_A(1+T)^D.
\]
Taking \(T=t/2\) in \eqref{eq:master-common-limit-bound} proves the
claimed rate for monomials, uniformly over \(\Mstar\). Every local
function is a finite linear combination of monomials.

Starting from \(\mu_{\mathbf1}\), compactness gives a subsequential
weak limit. Convergence of every monomial moment identifies all
subsequential limits with a single measure \(\pi\) satisfying
\eqref{eq:limit-moment-formula}. The Feller property gives
\(\pi P_s=\pi\).

If \(\nu\in\Mminus\), then
\(\bar\nu=(\nu+\mu_{\mathbf1})/2\in\Mstar\) by
\eqref{eq:Mminus-affine}, and
\(\nu=2\bar\nu-\mu_{\mathbf1}\). Linearity extends the estimate from
\(\Mstar\) to \(\Mminus\).
\end{proof}

\begin{proof}[Proof of \cref{cor:uniform-ergodicity}]
\Cref{lem:all-measures-Mminus} makes
\eqref{eq:common-limit-bound} uniform over all initial laws. If
\(\rho\) is invariant, then, for every local \(f\),
\[
  \abs{\rho(f)-\pi(f)}
  =
  \abs{\rho(P_tf)-\pi(f)}
  \le
  \sup_\xi\abs{P_tf(\xi)-\pi(f)}
  \longrightarrow0.
\]
Thus \(\rho=\pi\). Finally,
\((1+t)^De^{-\varepsilon t/2}\le C_{D,\varepsilon}
e^{-\varepsilon t/4}\), which gives
\uniformexponentialergodicity.
\end{proof}
